\chapter{Notation and Conventions}

\begin{table}[htbp]
\centering
\caption{Core mathematical, DSP, and RTOS symbols used throughout this book.}
\begin{tabular}{@{}lp{9.6cm}@{}}
\toprule
Symbol & Meaning \\
\midrule
$f_s$ & Audio sample rate in Hz; restricted to $\{44100, 48000\}$ (default $48000$). \\
$Q$ & Synthesis block quantum in frames; fixed at $Q = 32$ (\texttt{RENDER\_BLOCK}). \\
$T$ & Audio block period $T = Q/f_s$: $666.67\,\mu\text{s}$ at $48\,\text{kHz}$, $725.62\,\mu\text{s}$ at $44.1\,\text{kHz}$. \\
$H_{\text{per}}$ & Host VirtIO audio period; fixed at $1024$ frames ($21.33\,\text{ms}$ / $23.22\,\text{ms}$). \\
$\mathrm{WCET}$ & Declared worst-case execution time per quantum ($632\,\mu\text{s}$ for Orchestra). \\
$n \in [0, 127]$ & Standard MIDI pitch index; $n = 69$ corresponds to concert $\mathrm{A}_4 = 440\,\text{Hz}$. \\
$\mathbb{Z}_{12}$ & Cyclic group of 12-tone equal-tempered pitch classes modulo the octave. \\
$D_{12}$ & Dihedral group of order $24$ generated by transposition $T_k$ and inversion $I_k$. \\
$\mathcal{P}, \mathcal{L}, \mathcal{R}$ & Neo-Riemannian involutions (Parallel, Leittonwechsel, Relative) on triads. \\
$\mathrm{PPQ}$ & Pulses per quarter note; fixed at $\texttt{ORCH\_PPQ} = 960$ ticks. \\
$v_m \in [0, 127]$ & MIDI velocity mapped to amplitude $v \in [0.08, 0.45]$ via curves $v_0, v_1, v_2$. \\
$G_{\mathrm{trk}}, G_L, G_R$ & Track fader gain and constant-power stereo panning gains. \\
$H(z)$ & Discrete-time $z$-domain transfer function of a linear time-invariant filter. \\
$C_k(x)$ & Non-linear WaveShaper saturation curve with drive parameter $k \in [0, 100]$. \\
$J_m(I)$ & Bessel function of the first kind of order $m$ at FM modulation index $I$. \\
$S(t_0, \mathbf{p})$ & Stateless procedural Sound Shim triggered at $t_0$ with parameter vector $\mathbf{p}$. \\
$S^3 \to S^2$ & 4D Hopf fibration projected stereographically into $\mathbb{R}^3$. \\
$\zeta, \omega_i$ & Damping ratio ($\zeta = 1.0$, critical damping) and natural angular frequency. \\
\texttt{OrchView} & Plain-old-data (POD) desk snapshot published across a seqlock boundary. \\
\bottomrule
\end{tabular}
\end{table}

\noindent C99 identifiers, struct names, and function prototypes appear in \texttt{monospace}. All binary containers (\texttt{.orch}, \texttt{.mid}, \texttt{.wav}) specify explicit little-endian or big-endian byte layouts; compiler struct padding is never normative.
